% 05_multilingual_critical.tex
% Multilingual text and critical apparatus for Digital Humanities
% Overleaf compatible: YES (polyglossia + XeLaTeX/LuaLaTeX)

% !TeX program = xelatex
% !TeX encoding = UTF-8

\documentclass[12pt,a4paper]{article}
\usepackage{fontspec}
\usepackage{polyglossia}
\usepackage{hyperref}
\usepackage{cleveref}
\usepackage{reledmac}
\usepackage{reledpar}
\usepackage{parallel}
\usepackage{verse}
\usepackage{geometry}

% Language setup
\setmainlanguage{english}
\setotherlanguages{latin,greek,french,german,hebrew,arabic}

% Fonts for different scripts
\newfontfamily\latinfont{Latin Modern Roman}
\newfontfamily\greekfont{GFS Didot}
\newfontfamily\hebrewfont{David CLM}
\newfontfamily\arabicfont{Amiri}

% Critical apparatus setup
\Xarrangement[A]{paragraph}
\Xarrangement[B]{paragraph}
\Xwrapexplanations[A]{\par}
\Xwrapexplanations[B]{\par}
\XnotesWidth[A]{0.9\textwidth}
\XnotesWidth[B]{0.9\textwidth}
\linenumincrement*{5}
\firstlinenum*{1}

% Parallel text setup
\setlength{\ParallelLTextWidth}{0.45\textwidth}
\setlength{\ParallelRTextWidth}{0.45\textwidth}
\setlength{\ParallelAtEnd{\vspace{1em}}}

\hypersetup{
    colorlinks=true,
    linkcolor=blue,
    citecolor=red,
    urlcolor=cyan
}

\title{Multilingual Text and Critical Apparatus}
\author{Author Name}
\date{\today}

\begin{document}

\maketitle

\begin{abstract}
This example demonstrates multilingual typesetting and critical apparatus
for Digital Humanities: parallel texts, critical editions with apparatus
criticus, manuscript collation, and mixed-script documents.
\end{abstract}

\section{Multilingual Typesetting with Polyglossia}\label{sec:polyglossia}
Main text in English.

\begin{french}
Ceci est du texte en français.
\end{french}

\begin{german}
Dies ist ein deutscher Text.
\end{german}

\begin{latin}
Hic est textus latine scriptus.
\end{latin}

\begin{greek}
Αὐτὸς ἐστιν ὁ λόγος Ἑλληνιστὶ γεγραμμένος.
\end{greek}

\begin{hebrew}
זֶה טֶקְסְט בְּעִבְרִית.
\end{hebrew}

\begin{arabic}
هذا نص باللغة العربية.
\end{arabic}

Mixed inline: English, \textfrench{français}, \textgerman{Deutsch},
\textlatin{Latine}, \textgreek{Ἑλληνιστί}, \texthebrew{עִבְרִית}, \textarabic{العربية}.

\section{Right-to-Left Text}\label{sec:rtl}
Arabic paragraph:
\begin{arabic}
الفصل الأول من الكتاب يبدأ هنا. النص يكتب من اليمين إلى اليسار.
الأرقام تكتب من اليسار إلى اليمين: ١٢٣٤٥.
\end{arabic}

Hebrew paragraph:
\begin{hebrew}
הפרק הראשון של הספר מתחיל כאן. הטקסט נכתב מימין לשמאל.
המספרים נכתבים משמאל לימין: 12345.
\end{hebrew}

\section{Parallel Texts (Facing Translation)}\label{sec:parallel}
\begin{Parallel}{0.48\textwidth}{0.48\textwidth}
\ParallelLText{
\begin{latin}
\textbf{Latin Original:}\\
In principio creavit Deus caelum et terram.\\
Terra autem erat inanis et vacua,\\
et tenebrae super faciem abyssi.\\
Et spiritus Dei ferebatur super aquas.
}
\ParallelRText{
\textbf{English Translation:}\\
In the beginning God created the heaven and the earth.\\
And the earth was without form, and void;\\
and darkness was upon the face of the deep.\\
And the Spirit of God moved upon the face of the waters.
}
\ParallelPar
\end{Parallel}

\section{Critical Apparatus with reledmac}\label{sec:apparatus}
\begin{numbering}{section}
\beginnumbering
\pstart
\edtext{In principio}{\Afootnote{In initio \ms{A} \ms{B}; In principio \ms{C} \ms{D}}}
\edtext{creavit}{\Afootnote{creavit \ms{A} \ms{B} \ms{C}; fecit \ms{D}}}
Deus \edtext{caelum}{\Bfootnote{caelum \ms{A} \ms{B}; caelos \ms{C} \ms{D}}}
et terram.
\pend
\pstart
Terra autem \edtext{erat}{\Afootnote{erat \ms{A} \ms{B} \ms{C}; fuit \ms{D}}}
inanis et vacua,
et tenebrae \edtext{super}{\Afootnote{super \ms{A} \ms{B}; supra \ms{C} \ms{D}}}
faciem abyssi.
\pend
\endnumbering
\end{numbering}

\section{Apparatus with Multiple Layers (A=substantive, B=orthographic)}\label{sec:multiapparatus}
\begin{numbering}{section}
\beginnumbering
\pstart
\edtext{In principio}{\Afootnote{In initio \ms{A}; In principio \ms{B} \ms{C} \ms{D}}%
  \Bfootnote{principio] principio \ms{A} \ms{B}; principio \ms{C} \ms{D}}}
creavit Deus caelum et terram.
\pend
\endnumbering
\end{numbering}

\section{Manuscript Sigla and Witnesses}\label{sec:sigla}
% Define manuscript sigla
\def\msA{\textit{A}}  % Codex Vaticanus
\def\msB{\textit{B}}  % Codex Sinaiticus
\def\msC{\textit{C}}  % Codex Alexandrinus
\def\msD{\textit{D}}  % Codex Ephraemi

Critical edition with witness list:
\begin{itemize}
    \item \msA = Codex Vaticanus (4th c.)
    \item \msB = Codex Sinaiticus (4th c.)
    \item \msC = Codex Alexandrinus (5th c.)
    \item \msD = Codex Ephraemi Rescriptus (5th c.)
\end{itemize}

Variant reading:
\edtext{caelum}{\Afootnote{caelum \msA \msB; caelos \msC \msD}}

\section{Critical Edition Structure}\label{sec:structure}
A complete critical edition includes:

\begin{enumerate}
    \item \textbf{Introduction} -- History of the text, manuscripts, stemma
    \item \textbf{Text} -- Constituted text with line numbers
    \item \textbf{Apparatus Criticus} -- Substantive variants (A)
    \item \textbf{Apparatus Orthographicus} -- Orthographic variants (B)
    \item \textbf{Apparatus Fontium} -- Sources and parallels
    \item \textbf{Commentary} -- Explanatory notes
    \item \textbf{Indices} -- Verborum, nominum, locorum
\end{enumerate}

\section{Stemma Codicum}\label{sec:stemma}
% Using tikz-qtree or forest for stemma (see 09_tikz_graphics.tex)
% Simple text representation:
%
%                    [Archetype]
%                       |
%        +--------------+--------------+
%        |                             |
%     [α]                            [β]
%    /   \                          /   \
%  [A]  [B]                       [C]   [D]

\section{TEI XML Export Compatibility}\label{sec:tei}
% For TEI export, use structured markup:
% <div type="edition">
%   <p n="1">In principio <app><rdg wit="#A #B">In initio</rdg>
%   <rdg wit="#C #D">In principio</rdg></app> creavit Deus...</p>
% </div>
% <div type="apparatus">
%   <listApp>
%     <app n="1"><rdg wit="#A #B">In initio</rdg>
%     <rdg wit="#C #D">In principio</rdg></app>
%   </listApp>
% </div>

\end{document}